\chapter*{ESPECIFICACIONES TECNICAS}
\addcontentsline{toc}{chapter}{ESPECIFICACIONES TECNICAS}

% Reinicio de contadores para numeracion coherente con el capitulo 3.
\setcounter{section}{0}
\renewcommand{\thesection}{3.\arabic{section}.}
\renewcommand{\thesubsection}{3.\arabic{section}.\arabic{subsection}.}
\renewcommand{\theHsection}{especificaciones.\arabic{section}}
\renewcommand{\theHsubsection}{especificaciones.\arabic{section}.\arabic{subsection}}

\section{OBJETO, ALCANCE Y CRITERIO DE PRELACION}

Estas especificaciones establecen, para cada partida principal de las instalaciones
electricas interiores, los requisitos de materiales, el metodo de ejecucion, el control
de calidad y la forma de medicion y valorizacion. Se aplican a la vivienda unifamiliar
de dos pisos del proyecto Aquiles, con suministro monofasico de 220 V y 60 Hz.

Ante una diferencia entre documentos, se aplicara el siguiente orden de prelacion:
planos electricos aprobados, cuadro de cargas y diagrama unifilar, estas
especificaciones, metrados y presupuesto. Ninguna dimension o proteccion marcada como
preliminar podra emplearse para compra o construccion sin la revision del responsable
tecnico.

La base normativa principal es el Codigo Nacional de Electricidad -- Utilizacion
(CNE-U) y la Norma EM.010 Instalaciones Electricas Interiores del Reglamento Nacional
de Edificaciones. En particular, se consideran la Regla 030-002 para secciones minimas,
la Regla 020-132 para proteccion diferencial, la Regla 050-108 para reserva en tableros,
las Reglas 150-400 y 150-402 para proteccion y ubicacion de tableros, y la Regla
060-712 para resistencia de puesta a tierra.

\section{CANALIZACION EMPOTRADA EN TUBERIA PVC SAP}

\textbf{Unidad de medida:} metro (m).

\subsection{Requisitos de materiales y caracteristicas}
\begin{itemize}
  \item Tuberia rigida de PVC SAP para uso electrico, clase pesada, autoextinguible y
  resistente a humedad, impacto y aplastamiento durante la obra.
  \item Diametro nominal de 20 mm para circuitos derivados de 2.5 mm\(^2\). Los
  alimentadores y el circuito C6 emplearan el diametro que resulte del calculo de
  ocupacion y del escenario electrico aprobado; el modelo vigente propone 32 mm para
  conductores de 10 mm\(^2\) y 40 mm para el alimentador principal.
  \item Curvas de fabrica, uniones, conectores, pegamento para PVC y cajas compatibles.
  No se aceptaran tubos fisurados, deformados ni calentados directamente para formar
  curvas.
\end{itemize}

\subsection{Metodo de ejecucion}
\begin{enumerate}
  \item Replantear la ruta y las posiciones de cajas antes de picar. Verificar que la
  trayectoria no invada columnas, vigas, elementos estructurales, tuberias sanitarias
  ni otras instalaciones.
  \item Ejecutar el canaleteado con cortadora o amoladora provista de disco apropiado,
  cincel y martillo liviano. El ancho y la profundidad seran solo los necesarios para
  alojar la tuberia y su recubrimiento.
  \item Presentar, unir y fijar la tuberia sin aplastamientos. Entre dos cajas no se
  admitiran cambios de direccion que impidan el posterior tendido de conductores; de
  ser necesario se incorporara una caja de paso accesible.
  \item Taponar los extremos durante la obra civil, comprobar continuidad y ausencia
  de obstrucciones con una guia, y recien despues autorizar el resane.
\end{enumerate}

\subsection{Medicion, valorizacion y aceptacion}
Se medira la longitud neta instalada siguiendo el eje de la canalizacion, entre caras
de cajas, en metros con aproximacion a dos decimales. El precio unitario incluye
tuberia, accesorios ordinarios, adhesivo, fijacion, picado, resane basico, desperdicio,
mano de obra, herramientas, limpieza y prueba con guia. Las cajas y los accesorios
especiales se valorizaran en sus partidas correspondientes. Solo se aceptaran tramos
alineados, continuos, limpios y sin dano estructural.

\section{CONDUCTORES Y CABLEADO DE CIRCUITOS}

\textbf{Unidad de medida:} metro (m) de conductor instalado, identificado y probado.

\subsection{Requisitos de materiales y caracteristicas}
\begin{itemize}
  \item Conductores unipolares de cobre, con aislamiento certificado para 450/750 V o
  superior, aptos para canalizacion interior. Se priorizara aislamiento de baja emision
  de humos cuando el expediente definitivo asi lo especifique.
  \item Conforme a la Regla 030-002 del CNE-U, la seccion no sera menor de
  2.5 mm\(^2\) para circuitos derivados de alumbrado y fuerza, y de 1.5 mm\(^2\)
  solamente para control de alumbrado. El proyecto adopta 2.5 mm\(^2\) como minimo
  para los circuitos C1 a C5, C7 y C8.
  \item La seccion de C6 y de los alimentadores se definira con el escenario de carga,
  capacidad de corriente, caida de tension, proteccion y ocupacion de tuberia. El
  escenario conservador vigente propone 10 mm\(^2\) para C6 y para el alimentador de
  T2, sujeto a confirmacion de las cargas reales.
  \item Identificacion permanente de fase, neutro y conductor de proteccion PE; no se
  usara el color verde o verde/amarillo para otro fin.
\end{itemize}

\subsection{Metodo de ejecucion}
\begin{enumerate}
  \item Cablear solo cuando las tuberias esten terminadas, limpias, secas y verificadas
  con guia. No se permitiran empalmes dentro de la tuberia.
  \item Tender conjuntamente los conductores de cada circuito, evitando esfuerzos que
  danen el aislamiento. Usar lubricante compatible cuando sea necesario.
  \item Ejecutar derivaciones unicamente dentro de cajas accesibles, mediante bornes o
  conectores certificados. Dejar reserva suficiente y rotular ambos extremos.
  \item Verificar continuidad, identificacion, ausencia de cortocircuitos y resistencia
  de aislamiento antes de energizar.
\end{enumerate}

\subsection{Medicion, valorizacion y aceptacion}
Cada conductor se medira por su longitud real instalada, incluyendo reservas tecnicas
razonables en cajas y tableros. El precio incluye suministro, tendido, identificacion,
terminales ordinarios, mano de obra, herramientas y pruebas. Se rechazaran conductores
con aislamiento cortado, colores no conformes, empalmes ocultos o secciones distintas
de las aprobadas.

\section{SALIDAS DE ALUMBRADO E INTERRUPTORES}

\textbf{Unidad de medida:} punto (pto).

\subsection{Requisitos de materiales y caracteristicas}
\begin{itemize}
  \item Caja octogonal para centro de luz y caja rectangular para el dispositivo de
  mando, ambas adecuadas al sistema empotrado.
  \item Interruptor simple o conmutado de 10 A, 250 V, con placa aislante y bornes
  protegidos. Luminaria LED tipo plafon de 18 W como referencia de presupuesto.
  \item Conductores de 2.5 mm\(^2\) para alimentacion del circuito y conductor de
  control conforme al plano aprobado; toda masa metalica accesible debera enlazarse al
  conductor PE cuando corresponda.
\end{itemize}

\subsection{Metodo de ejecucion}
\begin{enumerate}
  \item Ubicar el centro de luz segun el plano y colocar el interruptor a 1.20 m sobre
  nivel de piso terminado, cerca del acceso y sin interferir con el giro de la puerta.
  \item Nivelar y enrasar las cajas con el acabado. Realizar conexiones dentro de las
  cajas, fijar la luminaria y proteger las partes activas.
  \item Probar continuidad, accionamiento simple o conmutado, encendido, apagado y
  correcta identificacion de fase y retorno.
\end{enumerate}

\subsection{Medicion, valorizacion y aceptacion}
Se medira por punto completo, instalado y operativo. La valorizacion de esta partida
incluye cajas, interruptor, placa, luminaria y accesorios de conexion; la canalizacion
y los conductores se valorizan por separado para evitar doble conteo. El punto se
aceptara nivelado, firmemente fijado y con prueba funcional satisfactoria.

\section{TOMACORRIENTES GENERALES 2P+T}

\textbf{Unidad de medida:} punto (pto).

\subsection{Requisitos de materiales y caracteristicas}
\begin{itemize}
  \item Tomacorriente doble 2P+T de 16 A, 250 V, con placa de material aislante
  autoextinguible y alveolos protegidos.
  \item Caja rectangular y terminal efectivo para el conductor de proteccion. El
  circuito empleara fase, neutro y PE de la seccion aprobada.
  \item En zonas expuestas a humedad o intemperie se usara envolvente con grado de
  proteccion adecuado a la ubicacion y proteccion diferencial de 30 mA.
\end{itemize}

\subsection{Metodo de ejecucion}
\begin{enumerate}
  \item Instalar los puntos generales a 0.30 m sobre el nivel de piso terminado, salvo
  indicacion distinta en planos. Los puntos sobre mesa se ubicaran aproximadamente a
  1.10 m y se coordinaran con muebles y aparatos.
  \item Conectar fase, neutro y PE sin invertir polaridad. Mantener la continuidad del
  PE aunque se retire el mecanismo para mantenimiento.
  \item Alinear la placa, ajustar bornes al torque indicado por el fabricante y probar
  polaridad, continuidad de proteccion y funcionamiento.
\end{enumerate}

\subsection{Medicion, valorizacion y aceptacion}
Se medira por punto terminado. El precio incluye caja, mecanismo, placa, terminales,
montaje y prueba; tuberia y conductores se pagan por separado. No se aceptaran placas
flojas, mecanismos hundidos, polaridad invertida ni ausencia de continuidad de tierra.

\section{SALIDAS ESPECIALES: COCINA, LAVADORA Y BOMBA}

\textbf{Unidad de medida:} punto (pto).

\subsection{Requisitos de materiales y caracteristicas}
\begin{itemize}
  \item Cada carga especial tendra circuito y proteccion acordes con su potencia de
  placa. La cantidad, corriente nominal, clavija, seccion de conductor y proteccion de
  C6, C7 y C8 deben cerrarse antes de compra.
  \item Para la bomba exterior se usara caja estanca de grado minimo IP55 o el exigido
  por el fabricante, prensaestopas y medio local de desconexion cuando corresponda.
  \item La proteccion diferencial tendra sensibilidad no mayor de 30 mA, sin sustituir
  al conductor de proteccion ni al sistema de puesta a tierra.
\end{itemize}

\subsection{Metodo de ejecucion}
\begin{enumerate}
  \item Confirmar potencia, corriente, regimen de servicio y manual del equipo antes de
  instalar el punto. No se conectaran varias cargas de alta potencia a un circuito sin
  recalcular demanda, conductor y proteccion.
  \item Coordinar altura y posicion con el equipo. En cocina y lavanderia se evitara la
  proximidad a fuentes de agua y calor; en exterior se sellaran todas las entradas.
  \item Rotular el circuito en el tablero y probar continuidad de PE, polaridad,
  aislamiento y accionamiento de la proteccion diferencial.
\end{enumerate}

\subsection{Medicion, valorizacion y aceptacion}
Se medira por salida especial instalada y probada. El precio incluye caja, mecanismo o
borne de conexion, sellos, accesorios y pruebas; canalizacion, conductores y proteccion
en tablero se valorizan en partidas separadas. Una salida no sera aceptada mientras la
carga de placa y su coordinacion de protecciones permanezcan sin confirmar.

\section{ALIMENTADORES Y TABLEROS TG/T2}

\textbf{Unidad de medida:} metro (m) para alimentadores y unidad (und) para tableros.

\subsection{Requisitos de materiales y caracteristicas}
\begin{itemize}
  \item Conductores, tuberia y proteccion seleccionados por capacidad de corriente,
  caida de tension, cortocircuito y coordinacion. El suministro del proyecto es
  monofasico; no se empleara la denominacion ``alimentador trifasico''.
  \item Gabinete con riel DIN, barras independientes de neutro y PE, puerta, directorio
  de circuitos y espacio para dos dispositivos de proteccion adicionales conforme a
  la Regla 050-108(2).
  \item Cada circuito derivado tendra interruptor termomagnetico bipolar. Los
  diferenciales seran de 30 mA y se agruparan como maximo de a tres circuitos, o se
  instalara uno por circuito, de acuerdo con la Regla 150-400. La corriente nominal de
  cada diferencial sera igual o mayor que la proteccion ubicada aguas arriba.
  \item La indicacion previa de tableros de 8 polos es insuficiente para cerrar la
  compra: el numero final de modulos debe obtenerse del diagrama unifilar considerando
  interruptor general, diferenciales, termomagneticos bipolares y reserva.
\end{itemize}

\subsection{Metodo de ejecucion}
\begin{enumerate}
  \item Instalar los tableros en lugar seco, accesible y fuera de baños, escaleras,
  armarios y ambientes de doble altura. Ninguna manija quedara a mas de 1.70 m sobre el
  piso, conforme a la Regla 150-402.
  \item Fijar el gabinete, montar equipos en riel DIN, separar neutro y PE, tender
  conductores ordenados, colocar terminales y ajustar bornes al torque del fabricante.
  \item Rotular tablero, alimentadores y circuitos C1 a C8. Colocar el directorio de
  circuitos en la contratapa y verificar la correspondencia con planos.
  \item Probar continuidad de PE, aislamiento, secuencia de maniobra y boton TEST de
  cada diferencial antes de energizar cargas.
\end{enumerate}

\subsection{Medicion, valorizacion y aceptacion}
El alimentador se medira por metro de cada conductor y por metro de tuberia. El tablero
se medira por unidad completa instalada y probada. La valorizacion incluye gabinete,
equipos indicados en el unifilar, barras, terminales, rotulado, conexionado y pruebas.
No se valorizara un tablero cuyo numero de modulos o coordinacion de protecciones no
este aprobado.

\section{SISTEMA DE PUESTA A TIERRA}

\textbf{Unidad de medida:} conjunto (cj).

\subsection{Requisitos de materiales y caracteristicas}
\begin{itemize}
  \item Electrodo tipo varilla de acero cobreado de 5/8 pulg o 3/4 pulg por 2.40 m,
  segun el diseño aprobado y disponibilidad certificada.
  \item Conductor de cobre de seccion calculada, caja de registro inspeccionable y
  conector certificado compatible con electrodo y conductor.
  \item El tratamiento de terreno, si se requiere, sera no corrosivo y se aplicara de
  acuerdo con su ficha tecnica. No reemplaza la medicion final.
\end{itemize}

\subsection{Metodo de ejecucion}
\begin{enumerate}
  \item Confirmar la ubicacion fuera de cimentaciones, redes enterradas y zonas donde
  la caja quede inaccesible. Excavar con dimensiones suficientes para instalar el
  electrodo vertical y ejecutar el tratamiento previsto.
  \item Instalar electrodo, conductor y conector sin uniones ocultas. Llevar el
  conductor hasta la barra PE del tablero y protegerlo contra daño mecanico.
  \item Colocar la caja de registro a nivel terminado, rellenar y compactar por capas.
  Medir la resistencia con telurimetro y registrar fecha, metodo, equipo y resultado.
\end{enumerate}

\subsection{Medicion, valorizacion y aceptacion}
Se medira por conjunto terminado e incluye excavacion, electrodo, conductor del tramo
considerado, conector, tratamiento, caja, relleno, limpieza y protocolo de medicion.
Conforme a la Regla 060-712, la resistencia no sera mayor de 25 \(\Omega\). Si un
electrodo simple supera ese valor, se corregira el sistema y se repetira la medicion
antes de su aceptacion.

\section{PRUEBAS ELECTRICAS Y PUESTA EN SERVICIO}

\textbf{Unidad de medida:} global (glb).

\subsection{Requisitos de equipos y registros}
\begin{itemize}
  \item Multimetro, megometro, telurimetro, comprobador de tomacorrientes y equipos de
  proteccion personal con calibracion o verificacion vigente cuando corresponda.
  \item Planos finales, cuadro de circuitos, formatos de prueba y etiquetas de
  observacion y conformidad.
\end{itemize}

\subsection{Metodo de ejecucion}
\begin{enumerate}
  \item Inspeccionar visualmente canalizaciones, cajas, tableros, apriete, rotulos y
  continuidad del conductor PE con la instalacion desenergizada.
  \item Medir resistencia de aislamiento entre conductores activos y respecto a tierra
  con el nivel de prueba apropiado para los equipos conectados.
  \item Comprobar polaridad, continuidad, funcionamiento de luminarias y
  tomacorrientes, disparo de diferenciales y resistencia del sistema de puesta a
  tierra.
  \item Registrar resultados, corregir no conformidades y repetir las pruebas
  afectadas. La energizacion final requiere acta de conformidad.
\end{enumerate}

\subsection{Medicion, valorizacion y aceptacion}
La partida se medira en forma global cuando se entregue el protocolo completo y todas
las observaciones esten cerradas. El precio comprende uso de instrumentos, personal,
formatos, rotulado, correcciones menores y limpieza final; no incluye reparaciones por
trabajos defectuosos, que seran responsabilidad de la partida que los origino.
